\subsection{SAR, Pipeline and Delta-Sigma DACs}\label{sec:dac_other}

\subsubsection{SAR (Serial) DAC}
\paragraph{Overview:} a small DAC reused once per bit: \(C_1\) is charged to \(b_k V_{REF}\) and shared
with an equal \(C_2\) holding the running result, LSB first.
\noindent\textbf{Schematic:}\par
\begin{circuitfig}
  \draw (0,2) node[circ]{} node[left]{$b_k V_{REF}$} to[nos, l={charge}] (2,2) node[circ]{} to[nos, l={share}] (4.5,2) node[circ]{} -- (6,2) node[circ]{} node[right]{$V^{(k)}$};
  \draw (2,2) to[capacitor, l_={$C_1$}] (2,0) node[ground]{};
  \draw (4.5,2) to[capacitor, l={$C_2$}] (4.5,0) node[ground]{};
\end{circuitfig}
\paragraph{Transfer Function:}
\begin{enumerate}
  \item \textbf{Large Signal, Time Domain, DC Input Analysis (Q5)}: charge conservation when \(C_1\)
    and \(C_2\) are joined:
\begin{align*}
  C_1 b_k V_{REF} + C_2 V^{(k-1)} &= (C_1 + C_2) V^{(k)} \\
  V^{(k)} &= \frac{C_1 b_k V_{REF} + C_2 V^{(k-1)}}{C_1 + C_2} = \frac{b_k V_{REF} + V^{(k-1)}}{2} \hspace{1cm} [C_1 = C_2,\ V^{(-1)} = 0] \\
  V^{(0)} &= \frac{b_0 V_{REF}}{2}, \hspace{0.5cm}
  V^{(1)} = \frac{b_1 V_{REF}}{2} + \frac{b_0 V_{REF}}{4}, \hspace{0.5cm} \ldots \hspace{0.5cm}
  V^{(k)} = V_{REF} \sum_{j=0}^{k} b_j 2^{j-(k+1)} \\
  V^{(n-1)} &= V_{REF} \sum_{k=0}^{n-1} b_k 2^{k-n}
\end{align*}
    the same recursion as the R-2R ladder (Section~\ref{sec:dac_r2r}), in time instead of space.
  \item \textbf{Small Signal, Time Domain, AC Input Analysis (Q2)}: the share switch \(R_{sw}\) between
    \(C_1\) and \(C_2\) (series \(\frac{C_1 C_2}{C_1 + C_2} = \frac{C}{2}\)): the difference \(V_{C1} - V_{C2}\) decays as
    \(e^{-t/T_c}\), \(T_c = \frac{R_{sw} C}{2}\).
\end{enumerate}
\paragraph{Analysis:}
\begin{enumerate}
  \item \textbf{DC Gain}: \(\frac{V^{(n-1)}}{D} = \frac{V_{REF}}{2^n} = 1\,\text{LSB}\).
  \item \textbf{Speed}: \(n\) charge/share cycles per conversion, each settled to \(\frac{1}{2^{n+1}}\):
    \begin{align*}
      t_{share} &= T_c \ln(2^{n+1}) = (n + 1)\ln(2)\, \frac{R_{sw} C}{2}
    \end{align*}
  \item \textbf{Mismatch}: \(C_1 = C(1 + e)\), \(C_2 = C\): the weight of each new bit is
    \begin{align*}
      \frac{C_1}{C_1 + C_2} &= \frac{1 + e}{2 + e} \approx \frac{1}{2}\left(1 + \frac{e}{2}\right) \hspace{1cm} [\tfrac{1}{1 + x} \approx 1 - x,\ e^2 \to 0]
    \end{align*}
    the ratio error \(\frac{e}{2}\) compounds over \(n\) bits: \(|e| \lesssim 2^{-n}\) for \(n\)-bit linearity.
  \item \textbf{Advantages \& Limitations}: two capacitors for any \(n\) (smallest area); but \(n\)
    cycles per conversion and a matching requirement of \(2^{-n}\).
\end{enumerate}

\subsubsection{Pipeline DAC}
\paragraph{Overview:} the serial recursion unrolled over \(n\) stages (one per bit, MSB last).
\noindent\textbf{Schematic:}\par
\begin{circuitfig}
  \foreach \i/\lab in {0/{$b_0$}, 3/{$b_1$}, 6.5/{$b_{n-1}$}} {
    \draw (\i,0) rectangle ++(2,1.2) node[midway, align=center]{$\frac{V + b V_{REF}}{2}$};
    \draw[->] (\i+1,2) node[above]{\lab} -- (\i+1,1.2);
  }
  \draw[->] (-1,0.6) node[left]{$0$} -- (0,0.6);
  \draw[->] (2,0.6) -- (3,0.6); \draw[->] (5,0.6) -- (5.5,0.6); \node at (5.95,0.6) {$\ldots$}; \draw[->] (6.1,0.6) -- (6.5,0.6);
  \draw[->] (8.5,0.6) -- (9.5,0.6) node[right]{$V^{(n-1)}$};
\end{circuitfig}
\paragraph{Transfer Function:}
\begin{enumerate}
  \item \textbf{Large Signal, Time Domain, DC Input Analysis (Q5)}: stage \(k\):
    \begin{align*}
      V^{(k)} &= \frac{V^{(k-1)} + b_k V_{REF}}{2} \hspace{1cm} \Rightarrow \hspace{1cm} V^{(n-1)} = V_{REF} \sum_{k} b_k 2^{k-n} \hspace{1cm} [\text{as the serial DAC}, V^{(-1)} = 0]
    \end{align*}
\end{enumerate}
\paragraph{Analysis:}
\begin{enumerate}
  \item \textbf{Speed}: a new code (its bits delayed one cycle per stage) enters every cycle:
    throughput \(f_s\), latency \(n\) cycles.
  \item \textbf{Advantages \& Limitations}: one conversion per cycle with simple stages; but \(n\)
    stages (area, power), latency \(n\), and each stage's \(\frac{1}{2}\) weight must match to \(2^{-n}\).
\end{enumerate}

\subsubsection{Delta-Sigma DAC}
\paragraph{Overview:} a digital modulator, a 1-bit DAC and an analog low pass.
\noindent\textbf{Schematic:}\par

\begin{circuitfig}
  \draw[->] (-1.2,0.6) node[left]{$D$ ($n$ bits)} -- (0,0.6);
  \draw (0,0) rectangle (2.4,1.2) node[midway, align=center]{digital\\modulator};
  \draw[->] (2.4,0.6) -- node[above]{\footnotesize 1 bit, $f_s$} (4,0.6);
  \draw (4,0) rectangle (5.6,1.2) node[midway, align=center]{1-bit\\DAC};
  \draw[->] (5.6,0.6) -- (6.6,0.6);
  \draw (6.6,0) rectangle (8.6,1.2) node[midway, align=center]{analog\\low pass};
  \draw[->] (8.6,0.6) -- (9.6,0.6) node[right]{$V_{OUT}$};
\end{circuitfig}
\paragraph{Transfer Function:}
\begin{enumerate}
  \item \textbf{Small Signal, Frequency Domain, AC Input Analysis (Q4)}: the digital loop of
    Section~\ref{sec:adc_ds} at \(f_s = 2\, OSR\, f_B\), \(Y = X + NTF\cdot E\) with \(NTF = (1 - z^{-1})^L\): the
    quantisation noise \(E\) is pushed above \(f_B\) and the analog low pass removes it.
\end{enumerate}
\paragraph{Analysis:}
\begin{enumerate}
  \item \textbf{Resolution}: for an \(L\)-th order modulator with an \(n_q\)-bit quantiser (derived for
    the Delta-Sigma ADC, Section~\ref{sec:adc_ds}):
    \[
      SQNR_{dB} = 6.02\, n_q + 1.76 - 10\log_{10}\frac{\pi^{2L}}{2L + 1} + (2L + 1)\, 10\log_{10}(OSR)
    \]
  \item \textbf{Advantages \& Limitations}: a 1-bit DAC is perfectly linear (two levels define a
    line), resolution from \(OSR\) and \(L\) (digital); but \(f_s = 2\,OSR\,f_B\) (low bandwidth), out of band
    noise the analog filter must remove, and clock jitter sensitivity at the 1-bit output.
\end{enumerate}
